Experiment~3 showed that the choice between the median/MAD and mean-based Stage-B estimators did not materially explain the performance of the proposed pipeline, evaluated separately at Fp1 and at Fp2, the two electrodes Experiment~1 identified as consistently best-performing, rather than at a best-channel-per-session oracle. The median/MAD centre was intended to reduce the influence of extreme values remaining in the suspicious-epoch pool, consistent with the robustness of median-based estimators to outlying observations~\citep{leys2013detecting,rousseeuw1993alternatives}. However, pooled precision fell from Proposed-Mean to Proposed-approach at both electrodes (80.49\% to 79.51\% at Fp1; 80.23\% to 79.31\% at Fp2), while recall rose (88.07\% to 89.83\% at Fp1; 86.86\% to 88.73\% at Fp2) and $F_1$ increased modestly (82.12\% to 82.62\% at Fp1; 81.54\% to 82.15\% at Fp2). The estimators were statistically indistinguishable at either electrode (two-tailed Wilcoxon on session-level $F_1$: $p=0.791$ at Fp1, $p=0.775$ at Fp2): pooled, Proposed-approach exceeded Proposed-Mean in 44/104 (Fp1) and 50/104 (Fp2) sessions, was equal in 2 and 4, respectively, and was lower in 58 and 50. The median centre therefore did not provide the hypothesised precision advantage and cannot be credited as a demonstrated source of improvement, at either electrode. It remains a defensible default on robustness grounds, but the measured comparison supports similar overall performance from the two estimators.
